\documentclass[12pt]{book}

\usepackage[acronym]{glossaries}
\usepackage{hyperref}
\usepackage{listings}
\usepackage{amsmath}
\usepackage{amssymb}
\usepackage{mathtools}
\usepackage{systeme}
\usepackage{algorithm,algpseudocode}
\usepackage{afterpage}
\usepackage[a4paper, total={6in, 10in}]{geometry}
\usepackage{scrextend}
\usepackage{tcolorbox}
\usepackage{polynom}

\pagestyle{plain}

\newcommand{\example}[1]{
    \begin{addmargin}[1em]{1em}
    \begin{tcolorbox}
    #1
    \end{tcolorbox}
    \end{addmargin}
}

\newcommand{\part}[2]{
    \renewcommand\chaptername{}
    \renewcommand\thechapter{}
    \chapter{#2}
    \renewcommand\thechapter{#1}
}

\newcommand{\code}[1]{
    {\color{blue}\texttt{#1}}
}

\definecolor{dkgreen}{rgb}{0,0.6,0}
\definecolor{gray}{rgb}{0.5,0.5,0.5}
\definecolor{mauve}{rgb}{0.58,0,0.82}

% This `if` statement stands here just to comment out 
% examples of figure and listing.
\iffalse
% Example of figure.
\begin{figure}[h]
\centering
\includegraphics[scale=0.5]{Resources/Lecture_2/co_np.png}
\caption{Co-NP problems}
\label{fig:co_np}
\end{figure}

% Example of listing.
\begin{lstlisting}[
    language=Verilog,
    caption=SystemVerilog packed/unpacked arrays,
    columns=flexible,
    basicstyle={\footnotesize\ttfamily},
    numbers=left,
    numberstyle=\tiny\color{gray},
    keywordstyle=\color{blue},
    commentstyle=\color{dkgreen},
    stringstyle=\color{mauve},
    tabsize=4
]
bit [3:0] data; // packed array
logic queue [3:0]; // unpacked array
\end{lstlisting}

% Example of table.
\begin{center}
\begin{tabular}{|c|c|c|}
    \hline
    & this entry & mirror entry \\
    \hline
    distance to the closest cluster & 10 & 6 \\
    \hline
    distance to the second closest cluster & 12 & 8 \\
    \hline
\end{tabular}
\end{center}
\fi



\begin{document}

\part{Čtení s porozuměním}
\section{Úloha 1}
Poslední zářijový víkend (1.pad, jednotne, hrad) bude patřit oslavám (3. pad, množné, žena)
svatého Václava (2. pad, jednonte, pán), patrona (2. pad, jednotne, pán) české země (2. pad, jednotne, růže), 
pivovarníků (2. pad, množné, pán) a vinařů (2. pad, množné, pán). 
Svatováclavské slavnosti (1. pad, množné, kost) nabídnou od pátku (2. pad, jednotne, hrad) 26. do 
neděle (2. pad, jednotne, hrad) 28. září (2. pad, jednotne, hrad) více než tři desítky 
zážitků (2. pad, jednotne, hrad).

\textbf{Speciální exkurze vás zavedou na běžně nepřístupná místa}
V sobotu (4. pad, jednotne, žena) večer (4. pad, jednotne, hrad) vyrazte na Svatováclavskou 
noc (4. pad, jednotne, kost) otevřených muzeí (2. pad, množné, muzeum) a galerií (2. pad, množné, muzeum). 
Vstup (1. pad, jednotne, muzeum) je zdarma! Mezi 20. a 24. hodinou (7. pad, jednonte, žena) si můžete 
prohlédnout třeba Městské divadlo (4. pad, jednotne, město), synagogu (2. pad, jednotne, žena), 
regionální muzeum (1. pad, jednotne, muzeum) nebo Museum Fotoateliér Seidel, které přichystalo 
program (4. pad, jednotne, hrad) také na sobotní a nedělní den (4. pad, jednotne, hrad). Při takzvaném 
Svatováclavském fotografování (6. pad, jednotne, stavení) zažijete atmosféru (4. pad, jednotne, žena) 
počátku (2. pad, jednotne, hrad) 20. století (2. pad, jednotne, stavení) a získáte historickou 
fotografii (2. pad, jednotne, růže) do svého rodinného alba (2. pad, jednotne, album). Vyrazit si pro ni 
můžete vždy mezi 9. a 15. hodinou (7. pad, jednotne, žena). V neděli (4. pad, jednotne, růže) si s 
průvodcem (7. pad, jednotne, soudce) můžete projít celý Český Krumlov (4. pad, jednotne, hrad). 
Vyráží se ve 14:30 z náměstí Svornosti (2. pad, jednotne, stavení). Prohlídka (1. pad, jednotne, žena) 
potrvá přibližně 90 minut (2. pad, množné, žena) a vstupenky (1. pad, množné, žena) zdarma jsou dostupné 
v Infocentru (6. pad, jednotne, centrum). Celý víkend se konají prohlídky (4. pad, množné, žena) 
zámku (2. pad, jednotne, hrad) (začátky (1. pad, množné, hrad) v 10:00, 11:00, 13:00, 14:00, 15:00, 16:00 
hodin (2. pad, množné, žena)). Vstupné (1. pad, jednotne, moře) pro dospělou osobu (4. pad, jednotne, žena) 
je 150 Kč, návštěvníci (1. pad, množné, pán) mladší 18 let a senioři (1. pad, množné, pán) mají 
vstup (4. pad, jednotne, hrad) zdarma. Vstupenky (1. pad, množné, žena) se prodávají v 
pokladně (6. pad, jednotne, žena) na nádvoří (1. pad, jednotne, stavení) zámku (2. pad, jednotne, hrad).

\textbf{Grafitový důl}
V neděli (4. pad, jednotne, růže) se zpřístupní také podzemí (4. pad, jednotne, stavení) 
grafitového dolu (2. pad, jednotne, hrad). Místo (2. pad, jednotne, město) můžete prozkoumat za jízdy 
na důlním vláčku (6. pad, jednotne, hrad) s průvodci (7. pad, jednotne, soudce). Před 
cestou (7. pad, jednotne, žena) vám (3. pad, množné, vy) zapůjčí ochranné pomůcky (4. pad, množné, žena), 
hornický oděv (4. pad, jednotne, hrad) a obuv (4. pad, jednotne, kost), bez kterých vás na 
prohlídku (4. pad, jednotne, žena) nepustí. Prohlídky (1. pad, množne, žena) začínají v 
neděli (4. pad, jednotne, růže) ve 12, 13 a 14 hodin (2. pad, množné, žena) a trvají 80 
minut (2. pad, množné, žena).

\textbf{Letohrádek a zámecké sklepy}
Netradiční podívanou (7. pad, jednotne, žena) nabídne v neděli (4. pad, jednotne, růže) 
podzemí (6. pad, jednotne, stavení) letohrádku (2. pad, jednotne, hrad) Bellarie v zámecké 
zahradě (6. pad, jednotne, žena). Od 10 do 14 hodin (2. pad, množné, žena) budete mít 
příležitost (4. pad, jednotne, kost) nahlédnout do umělé jeskyně (2. pad, jednotne, růže). 
V přízemí (6. pad, jednotne, stavení) letohrádku (2. pad, jednotne, hrad) vás čeká 
prohlídka (1. pad, jednotne, žena) kuchyně (2. pad, jednotne, růže) a místnosti (2. pad, jednotne, kost) 
s výtahem (7. pad, jednotne, hrad) na jídlo (4. pad, jednotne, město). Ve stejném 
čase (6. pad, jednotne, hrad) budete moci vstoupit také do zámeckých 
sklepů (2. pad, jednotne, hrad) na III. nádvoří (4. pad, jednotne, stavení). 
Průvodci (1. pad, množné, soudce) budou vyrážet každých 30 minut, v podzemí (6. pad, jednotne, stavení) 
se zdržíte asi hodinu (4. pad, jednotne, žena). A hned poté můžete za hudbou (7. pad, jednotne, žena), 
a to do centra (2. pad, jednotne, centrum) na náměstí Svornosti (4. pad, jednotne, stavení).

\textbf{Centrum města rozezní koncerty}
Páteční úvod (1. pad, jednotne, hrad) hudebního festivalu (2. pad, jednotne, hrad) roztančí v 17:15 
českokrumlovská swingová kapela (1.pad, jednotne, žena) Black Bottom se světovou a českou 
klasikou (7. pad, jednotne, žena) počátku 20. století (2. pad, jednotne, stavení). Po nich v 19 
hodin (2. pad, množne, žena) nastoupí česká rocková legenda (1. pad, jednotne, žena), 
zpěvák (1. pad, jednotne, pán), skladatel (1. pad, jednotne, pán) a 
klávesista (1. pad, jednotne, předseda) Roman Dragoun. České rockové a popové 
hity (1. pad, množné, hrad) rozezní náměstí (4. pad, jednotne, stavení) 
Svornosti (2. pad, jednotne, kost) v sobotu (4. pad, jednotne, žena) od 20:30, kdy vystoupí na 
jednom pódiu (6. pad, jednotne, muzeum) mladá kapela (1. pad, jednotne, žena) z Českého Krumlova 
Chlapi v sobě a úspěšný českokrumlovský dětský sbor (1. pad, jednotne, hrad) Medvíďata s českými 
a moravskými lidovými písněmi (7. pad, množné, písen').



\end{document}
